\subsection{Rationalizing the Denominator (basic)} 
When we take a root of a fraction, we end up with a radical in the denominator according to the following property of roots, which corresponds to a property of exponents.
\begin{theorem}[Root of a Quotient Rule]
For any real numbers $x$ and $y\neq 0$ and any positive integer $p$ where all the roots below are real numbers,

\nic$\sqrt[p]{\dfrac{x}{y}}=\dfrac{\sqrt[p]{x}}{\sqrt[p]{y}}$
\end{theorem}
\begin{example}Simplify $\sqrt[3]{\dfrac{10}{27}}$.
\begin{lpic}{}{2}
\fractextleft{0.25}{-1}{$\sqrt[3]{\dfrac{10}{27}}={}$};
\end{lpic}
\end{example}
\noi We consider a quotient of radical expressions to be simplified only when there are no radicals in the denominator (for reasons that will become clear later). If we have a fraction with one term in the denominator, we use the following procedure.
\begin{procedure}{Rationalizing the Denominator (basic)}
To rationalize the denominator of a fraction $\dfrac{X}{\sqrt[p]{a}}$:
\begin{enumerate*}
\item Find the smallest number $b$ such that $ab$ is a perfect $p$th power. (A table on a calculator may be helpful here.)
\item Multiply the numerator and denominator by $\sqrt[p]{b}$.
\item Simplify the numerator and denominator, combining radicals if possible.
\end{enumerate*}
\end{procedure}
\begin{example}Simplify $\sqrt[4]{\dfrac{7}{8}}$.
\begin{lpic}{}{5}
\fractextleft{0.25}{-1}{$\sqrt[4]{\dfrac{7}{8}}={}$};
\regtextleft{0.25}{-3}{Trying 4th powers...};
\end{lpic}
\end{example}